\documentclass[11pt]{article}

\usepackage[T1]{fontenc}
\usepackage[utf8]{inputenc}
\usepackage{lmodern}
\usepackage{amsmath,amssymb,amsthm,mathtools}
\usepackage[a4paper,margin=28mm]{geometry}
\usepackage{microtype}
\usepackage{xurl}
\usepackage[hidelinks]{hyperref}
\usepackage{enumitem}
\usepackage{booktabs}
\usepackage{array}

\newtheorem{theorem}{Theorem}[section]
\newtheorem{lemma}[theorem]{Lemma}
\newtheorem{proposition}[theorem]{Proposition}
\newtheorem{corollary}[theorem]{Corollary}
\newtheorem{definition}[theorem]{Definition}
\newtheorem{remark}[theorem]{Remark}

\newcommand{\N}{\mathbb N}
\newcommand{\WIN}{\ensuremath{\mathrm{WIN}}}
\newcommand{\LOSS}{\ensuremath{\mathrm{LOSS}}}
\newcommand{\DRAW}{\ensuremath{\mathrm{DRAW}}}
\newcommand{\odd}{\operatorname{odd}}
\newcommand{\src}{\operatorname{src}}
\newcommand{\ct}{\operatorname{ct}}
\newcommand{\asl}{\operatorname{asl}}
\newcommand{\Moves}{\operatorname{moves}}

\title{The Two-Player $3n\mathbin{\pm}1$ Game:\
A Binary Normal Form and a Well-Founded No-Draw Proof}
\author{Grisha Pochuev\\
  \small\href{mailto:n_854@mail.ru}{n\_854@mail.ru}\\
  \small\url{https://github.com/Grisha-Pochuev/3n-plus-minus-1-game}}
\date{10 August 2026}

\begin{document}
\maketitle

\begin{abstract}
Consider the impartial two-player game on positive odd integers in which a
move from $n>1$ chooses $3n-1$ or $3n+1$ and removes all powers of two; reaching
$1$ wins immediately. Arbitrary play need not terminate: for example
$5\to 7\to 5$ is a legal cycle. The question is whether optimal play nevertheless
has finite remoteness from every starting position. We conjugate the game to a
binary normal form with an expanding move $A$ and a strictly contracting move
$B$, introduce constant-tail coordinates and a three-free source map, and rank
marked proof configurations by retained source, finite proof-height tokens, and
a typed fixed-fibre normalizer. The key entry step is a factor-removal lemma
that either removes one factor of three at the same tail exponent or produces
finite outcome provenance. This prevents an unranked return from a factor-free
three/two frame to the preceding exponent-one state. A one-shot initialization
then allows the typed normalizer to be entered without treating a larger
temporary inner source as a source decrease. The resulting well-founded marked
normalization excludes all draw positions.
\end{abstract}

\tableofcontents

\section{The game and the theorem}

The current state is a positive odd integer $n$. For $n>1$ the player to move
chooses one of
\[
  3n-1,\qquad 3n+1,
\]
and divides the chosen value by the largest possible power of $2$, leaving an
odd integer. If the resulting odd integer is $1$, the player who made the move
wins immediately. Players alternate with perfect information.

A position is called \WIN\ for the player to move if at least one move leads to
\LOSS. It is called \LOSS\ if every legal move leads to \WIN. These two classes
are the inductively generated finite-remoteness classes. Any position not in
either class is called \DRAW.

\begin{theorem}[Main theorem]\label{thm:main}
Every positive odd starting position has finite remoteness. Equivalently, the
$3n\pm1$ game has no \DRAW\ positions, and optimal play reaches $1$ after
finitely many moves.
\end{theorem}

The theorem is not a statement about arbitrary play. There are legal cycles,
such as
\[
  5\longrightarrow 7\longrightarrow 5,
\]
under suitable choices by the players.

\section{Binary conjugation}

Write an odd state as $n=2m+1$ and define
\[
  F(m)=\left\lceil\frac{3m}{2}\right\rceil .
\]
For a nonnegative integer $x$, let $R(x)$ be the integer obtained by deleting
the maximal alternating suffix of the binary expansion of $x$. Thus
$R(1001_2)=10_2$ and $R(10101_2)=0$.

\begin{proposition}[First normal form]\label{prop:first-normal}
For $m>0$, the two children in $m$-coordinates are exactly
\[
  F(m),\qquad F(R(m)).
\]
\end{proposition}

\begin{proof}
One has
\[
  3(2m+1)-1=2(3m+1),\qquad
  3(2m+1)+1=2(3m+2).
\]
Exactly one of $3m+1$ and $3m+2$ is odd. In the other expression, each
additional factor of two removed corresponds to crossing one additional
alternating bit at the end of the binary expansion of $m$. The process stops
at the first repeated adjacent bit, which is precisely the deletion defining
$R(m)$. Returning to $m$-coordinates gives the displayed pair.
\end{proof}

Since both children lie in the image of $F$, we represent the state $F(q)$ by
$q$. The conjugated game has terminal state $0$ and moves
\[
  A(q)=\left\lceil\frac{3q}{2}\right\rceil,
  \qquad
  B(q)=R(A(q)).
\]

\begin{proposition}[Strict contraction]\label{prop:B-contracts}
For every $q>0$,
\[
  B(q)<q.
\]
\end{proposition}

\begin{proof}
Deleting a nonempty binary suffix gives
\[
  B(q)\le \left\lfloor\frac{A(q)}2\right\rfloor
       \le \left\lfloor\frac{3q+1}{4}\right\rfloor<q
\]
for $q\ge2$, while $B(1)=0$.
\end{proof}

Thus all unbounded difficulty comes from the expanding branch $A$ and from the
unbounded alternating suffix removed by $B$.

\section{Constant-tail coordinates and coefficient sources}

For a positive odd integer $a$, an exponent $r\ge1$, and a phase
$e\in\{0,1\}$, put
\[
  Q_r^e(a)=a2^r-e.
\]
This is a binary word whose final constant block has length $r$ and consists
of the bit $e$.

\begin{lemma}[Long-tail recurrence]\label{lem:long-tail}
For every positive odd $a$, phase $e$, and $r\ge3$,
\[
  A(Q_r^e(a))=Q_{r-1}^e(3a),\qquad
  B(Q_r^e(a))=Q_{r-2}^e(3a).
\]
Moreover
\[
  B(Q_1^e(a))=B(Q_2^e(a)).
\]
\end{lemma}

\begin{proof}
For $r\ge3$, the expanding child still ends in at least two equal bits; hence
its maximal alternating suffix consists only of the final bit. This gives both
long-tail identities. At the boundary, the binary words of $3a-e$ and $6a-e$
differ by appending one bit that extends the same alternating suffix, so the
same prefix remains after deletion.
\end{proof}

Define the embedded three-free coefficient
\[
  J(s)=2A(s)+1.
\]
Every positive odd coefficient $a$ has a unique factorization
\[
  a=3^kJ(s),\qquad k\ge0.
\]
The integer $s$ is called the coefficient source of $a$ and is denoted
$\src(a)$.

\begin{lemma}[Source boundary]\label{lem:source-boundary}
For every $s>0$ there is a phase
\[
  \alpha(s)=1-\left(\left\lfloor\frac{s}{2}\right\rfloor\bmod2\right)
\]
such that
\[
  \odd(3J(s)+1-2\alpha(s))=J(A(s))
\]
with $2$-adic valuation one, while the opposite phase selects $J(B(s))$ with
valuation at least two.
\end{lemma}

Multiplying a constant-tail coefficient by $3$ preserves its coefficient
source. Passing through a factor-free signed boundary with coefficient $J(s)$
therefore exposes exactly one ordinary $A/B$ move on the source.

\section{Finite proof tokens}

The terminal state $0$ is \LOSS\ with height $h(0)=0$. A finite \WIN\ state has
height
\[
  h(x)=1+\min\{h(y):y\text{ is a \LOSS\ child of }x\},
\]
and a finite \LOSS\ state has height
\[
  h(x)=1+\max\{h(y):y\text{ is a child of }x\}.
\]
Only finite states whose outcome and height comparison have actually been
established are inserted into the marked configuration. They are called proof
tokens.

\begin{lemma}[Multiset token rank]\label{lem:token-rank}
For a finite multiset $M$ of proof heights define
\[
  \Phi(M)=\mathop{\#}_{h\in M}\omega^h,
\]
the natural ordinal sum of the monomials $\omega^h$. Replacing one token of
height $H$ by any finite multiset of certified descendants of heights strictly
below $H$ strictly decreases $\Phi$.
\end{lemma}

\begin{proof}
In Cantor normal form the coefficient of $\omega^H$ decreases by one, and all
inserted terms have lower exponent. No finite collection of lower terms can
compensate for the loss at exponent $H$.
\end{proof}

\section{Factor removal at fixed tail exponent}

The next lemma is the entry device that prevents a factorful exponent-one
state from creating an unranked restart.

\begin{lemma}[Fixed-tail factor removal]\label{lem:fixed-tail-factor}
Let $c$ be positive and odd, $e\in\{0,1\}$, and $r\ge2$. Put
\[
  X=Q_r^e(3c),\qquad G=Q_r^e(c),\qquad H=Q_{r+1}^e(c).
\]
If $X$ is \DRAW, then either $G$ is \DRAW, or the local configuration supplies
finite outcome provenance suitable for a marked typed entry. More explicitly,
if $Y=A(G)$, then
\[
  \Moves(H)=\{X,Y\},
\]
and, whenever $G$ is not \DRAW, at least one finite \WIN/\LOSS\ token is forced
among $G,Y$ and the other child of $G$.
\end{lemma}

\begin{proof}
Since $r+1\ge3$, Lemma~\ref{lem:long-tail} gives
\[
  A(H)=Q_r^e(3c)=X,
  \qquad
  B(H)=Q_{r-1}^e(3c).
\]
For $r\ge3$ this second state is $A(G)$ by the same long-tail recurrence. For
$r=2$, $G=Q_2^e(c)$ and $A(G)=Q_1^e(3c)$, again equal to $B(H)$. Thus
$\Moves(H)=\{X,Y\}$.

Because $X$ is \DRAW, the state $H$ cannot be \LOSS. If $H$ is \WIN, then its
other child $Y$ is \LOSS, giving a finite token. Assume therefore that $H$ is
\DRAW. Then $Y$ cannot be \LOSS. If $G$ is \DRAW, we are in the first
alternative.

It remains to suppose that $G$ is finite. If $G$ is \LOSS, then all its
children, including $Y$, are \WIN. If $G$ is \WIN, then either $Y$ is a \LOSS\
witness or, since $Y$ is not \LOSS, the other child of $G$ is a \LOSS\ witness.
In all cases the displayed local diamond provides an actual finite outcome
token with proved incidence.
\end{proof}

\begin{corollary}[Removing all factor powers at exponent two]\label{cor:remove-to-base}
Let $a=J(s)$ and $k\ge1$. If
\[
  Q_2^e(3^ka)
\]
is \DRAW, then either a finite provenance token is obtained, or
\[
  Q_2^e(a)
\]
is \DRAW.
\end{corollary}

\begin{proof}
Apply Lemma~\ref{lem:fixed-tail-factor} repeatedly with $r=2$ and
$c=3^{i-1}a$. If the token alternative never occurs, the factor exponent is
reduced by one at each step until it reaches zero.
\end{proof}

\begin{lemma}[Exponent-one lift to exponent two]\label{lem:exp-one-lift}
Let $a$ be positive and odd and $k\ge1$. If
\[
  R_k=Q_1^e(3^ka)
\]
is \DRAW, then either
\[
  P_{k-1}=Q_2^e(3^{k-1}a)
\]
is \DRAW, or $P_{k-1}$ is \WIN\ and its other child is a finite \LOSS\ token.
\end{lemma}

\begin{proof}
The long-tail boundary identities give
\[
  A(P_{k-1})=R_k.
\]
Since $P_{k-1}$ has a \DRAW\ child, it cannot be \LOSS. If it is not \DRAW,
it is \WIN; its \DRAW\ child $R_k$ cannot witness that it is \WIN, so its other
child is \LOSS.
\end{proof}

\begin{proposition}[Arbitrary factorful exponent-one entry]\label{prop:factorful-entry}
Let $a=J(s)$ and $k>0$. If
\[
  Q_1^e(3^ka)
\]
is \DRAW, then after finitely many outcome-compatible local steps either a
finite provenance token has been produced, or
\[
  Q_2^e(a)
\]
is \DRAW.
\end{proposition}

\begin{proof}
Apply Lemma~\ref{lem:exp-one-lift}. If it produces a finite token, stop. If it
produces the exponent-two \DRAW\ state $Q_2^e(3^{k-1}a)$, apply
Corollary~\ref{cor:remove-to-base}. This either produces a finite provenance
token or removes all remaining factors of three at the same exponent two.
\end{proof}

\section{Factor-free exponent-two base entry}

It remains to analyze the base state produced by
Proposition~\ref{prop:factorful-entry}. Let $x>0$, put
\[
  a=J(x),\qquad P=Q_1^e(a),\qquad U=Q_2^e(a),
\]
\[
  b=B(P)=B(U),\qquad F=A(U)=Q_1^e(3a),\qquad g=1-e.
\]

Suppose $U$ is \DRAW\ and $x$ is the globally least coefficient source of a
\DRAW. The common child $b$ has canonical coefficient strictly below $J(x)$.
Since $a=J(x)$ is factor-free, this implies that the coefficient source of $b$
is below $x$. Therefore $b$ is not \DRAW. As a child of a \DRAW\ state it also
cannot be \LOSS; hence
\[
  b\text{ is \WIN},\qquad F\text{ is \DRAW}.
\]
The exact first-factor identity is
\[
  \boxed{9J(x)+1-2e=2^vJ(b)}. \tag{1}
\]
Consequently the signed child of $F$ is
\[
  A(F)=Q_v^g(J(b)).
\]

\begin{proposition}[Base-entry attachment]\label{prop:base-entry}
Every factor-free exponent-two \DRAW\ at a globally minimal coefficient source
has a finite continuation that either strictly lowers the retained source,
